\section{Related Work}\label{sec:related}

Horner and Goldberg \cite{horner1991} published the first use of a genetic algorithm (GA) for music composition. Since then, several directions of GA application to melody composition have emerged. One notable example is GenJam \cite{biles1997}, which uses an interactive genetic algorithm, in other words, a human-in-the-loop fitness evaluation. The system evolves two populations of musical elements (measures and phrases), with notes mapped to the scale of the current chord. The fitness of each individual comes from a human mentor, who gives real-time feedback by typing ``g'' for good or ``b'' for bad. This approach exposes the ``fitness bottleneck'': human evaluation limits the speed and scalability of the evolutionary process. In contrast, our work replaces the human mentor with jazz theory rules, so the fitness is evaluated automatically.

Besides interactive approaches, another line of research computes the fitness automatically from general music theory rules. Ozcan and Ercal \cite{ozcan2008} present AMUSE, a GA-based system that generates improvised melodies over a given harmonic context. Each gene is a scale degree, where 0 is a rest and the maximum value is a hold. Its fitness function combines ten objectives based on music theory, including rules for chord tones, relationships between notes, directions of notes, the beginning and ending notes of the melody, and the proportion of rests. In a Turing test, listeners had to tell apart solos generated by the GA from solos written by an amateur musician, and 52\% of them gave the wrong answer, which suggests that the generated solos are comparable to an amateur's. Our work builds on this approach, using a similar rest/hold encoding and a fitness that combines several objectives.

Mati\'c \cite{matic2010} also uses an automatic fitness, to compose short melodies of four bars. The genome follows the same rest/hold encoding, with one integer per eighth note, where 0 is a rest, the values from 1 to $n$ are pitches relative to a reference note, and $n+1$ is a hold. The fitness is minimized and computed bar by bar as a weighted sum of the quality of each interval, the mean and variance of the intervals compared with reference values, and a penalty for notes outside the key. Since the melodies are short, the GA uses no crossover, only three mutations: moving a note by an octave, changing one note, and swapping two neighboring notes. The author reports that the GA tends to produce many rests, because rests avoid bad intervals. In contrast, our work keeps crossover, at bar lines, and rests come from an explicit breathing rule instead of being a side effect of avoiding bad intervals.

Besides these general techniques, Wiggins et al. \cite{wiggins1998} argue that the results of a GA depend heavily on the domain-specific knowledge in its operators and fitness function. Their jazz solo generator represents a solo as a list of (scale degree, duration) genes over a chord progression, and its fitness checks eight features, such as illegal jumps, long notes, speed, and pattern matching. They found that too much weight on pattern matching led to repetitive and boring solos. They also note that their systems have neither a clear plan nor intention, and that large-scale structure cannot be expected from a GA with little domain knowledge, a limitation that our work shares.
